% GENERATED FILE. Edit canonical lessons, then run npm run generate:books.
% canonical-source-hash: fnv1a64:3bf751225901a8db
% canonical-lessons: GE-C308-anstossen, GE-C308-sekt, GE-C308-jubeln, GE-C308-beifall, GE-C308-applaudieren

\chapter{Prost!}
\label{ch:ge-prost}
\hlchaptermodality{308}

\begin{chapteropening}
\textbf{By the end of this chapter:} \emph{I can describe a happy moment: to clink glasses, sparkling wine, to cheer, applause and to applaud.}

Complete the last lesson of chapter 308: I can describe a happy moment: to clink glasses, sparkling wine, to cheer, applause and to applaud.
\end{chapteropening}

\section[anstoßen]{anstoßen — to clink glasses, to drink to}

\label{lesson:GE-C308-anstossen}



\begin{quote}
Before the new one: say the German for the speech, then the German for the toast.
\end{quote}

\subsection*{You'll want to know: anstoßen}
\textbf{anstoßen} — \textquotedblleft{}to clink glasses, to drink to\textquotedblright{}.

\begin{quote}
Um Mitternacht stießen alle an.
\end{quote}

\begin{quote}
\emph{At midnight everyone clinked glasses.}
\end{quote}

Separable: \emph{wir stoßen an}. \emph{Auf jemanden anstoßen} — to drink to someone.

\begin{cousinweb}
\emph{anstoßen} is built on \emph{stoßen} \textquotedblleft{}to push, to bump\textquotedblright{}, from Proto-Germanic \emph{stautaną} (EtymDB), with \emph{an}: glass knocks on glass.
\end{cousinweb}

\begin{grammarlens}[title={What you've built}]
One more word for telling, explaining and describing at B1.
\end{grammarlens}

\subsection*{Your turn}
\begin{itemize}
\raggedright
  \item \emph{Say it:} \emph{anstoßen}
  \item \emph{Say it:} \emph{anstoßen}, once more
  \item \emph{Say it:} a sentence of your own with \emph{anstoßen}
  \item \emph{Recall:} say the German for the speech, then the German for the toast, then say \emph{anstoßen} again
\end{itemize}

\begin{tcolorbox}[breakable,colback=teal!4,colframe=teal!35!black,title={Before you move on}]
What does anstoßen mean? (\textquotedblleft{}To clink glasses, to drink to\textquotedblright{}.) Say it once more.
\end{tcolorbox}
\section[der Sekt]{der Sekt — the sparkling wine}

\label{lesson:GE-C308-sekt}



\begin{quote}
Before the new one: say the German for the toast, then the German for to clink glasses.
\end{quote}

\subsection*{You'll want to know: der Sekt}
\textbf{der Sekt} — \textquotedblleft{}the sparkling wine\textquotedblright{}.

\begin{quote}
Zur Feier gab es eine Flasche Sekt.
\end{quote}

\begin{quote}
\emph{For the celebration there was a bottle of sparkling wine.}
\end{quote}

German sparkling wine: \emph{ein Glas Sekt}.

\begin{cousinweb}
\emph{Sekt} is usually traced to French \emph{vin sec} or Italian \emph{vino secco} \textquotedblleft{}dry wine\textquotedblright{}, from Latin \emph{siccus} \textquotedblleft{}dry\textquotedblright{} (Lewis \& Short); EtymDB lists no parent.
\end{cousinweb}

\begin{grammarlens}[title={What you've built}]
One more word for telling, explaining and describing at B1.
\end{grammarlens}

\subsection*{Your turn}
\begin{itemize}
\raggedright
  \item \emph{Say it:} \emph{der Sekt}
  \item \emph{Say it:} \emph{der Sekt}, once more
  \item \emph{Say it:} a sentence of your own with \emph{der Sekt}
  \item \emph{Recall:} say the German for the toast, then the German for to clink glasses, then say \emph{der Sekt} again
\end{itemize}

\begin{tcolorbox}[breakable,colback=teal!4,colframe=teal!35!black,title={Before you move on}]
What does der Sekt mean? (\textquotedblleft{}The sparkling wine\textquotedblright{}.) Say it once more.
\end{tcolorbox}
\section[jubeln]{jubeln — to cheer, to rejoice}

\label{lesson:GE-C308-jubeln}



\begin{quote}
Before the new one: say the German for to clink glasses, then the German for the sparkling wine.
\end{quote}

\subsection*{You'll want to know: jubeln}
\textbf{jubeln} — \textquotedblleft{}to cheer, to rejoice\textquotedblright{}.

\begin{quote}
Die Leute jubelten, als die Mannschaft gewann.
\end{quote}

\begin{quote}
\emph{People cheered when the team won.}
\end{quote}

Regular: \emph{jubelte}, \emph{hat gejubelt}.

\begin{cousinweb}
\emph{jubeln} is built on \emph{Jubel} \textquotedblleft{}rejoicing\textquotedblright{}, from Middle High German \emph{jūbel}, which came from Latin \emph{iūbilus} \textquotedblleft{}a shout of joy\textquotedblright{} (EtymDB).
\end{cousinweb}

\begin{grammarlens}[title={What you've built}]
One more word for telling, explaining and describing at B1.
\end{grammarlens}

\subsection*{Your turn}
\begin{itemize}
\raggedright
  \item \emph{Say it:} \emph{jubeln}
  \item \emph{Say it:} \emph{jubeln}, once more
  \item \emph{Say it:} a sentence of your own with \emph{jubeln}
  \item \emph{Recall:} say the German for to clink glasses, then the German for the sparkling wine, then say \emph{jubeln} again
\end{itemize}

\begin{tcolorbox}[breakable,colback=teal!4,colframe=teal!35!black,title={Before you move on}]
What does jubeln mean? (\textquotedblleft{}To cheer, to rejoice\textquotedblright{}.) Say it once more.
\end{tcolorbox}
\section[der Beifall]{der Beifall — the applause}

\label{lesson:GE-C308-beifall}



\begin{quote}
Before the new one: say the German for the sparkling wine, then the German for to cheer.
\end{quote}

\subsection*{You'll want to know: der Beifall}
\textbf{der Beifall} — \textquotedblleft{}the applause\textquotedblright{}.

\begin{quote}
Am Ende gab es langen Beifall.
\end{quote}

\begin{quote}
\emph{At the end there was long applause.}
\end{quote}

No plural. \emph{Viel Beifall bekommen} — to get a lot of applause.

\begin{cousinweb}
\emph{Beifall} joins \emph{bei} \textquotedblleft{}with\textquotedblright{} and \emph{Fall}, from \emph{fallen} \textquotedblleft{}to fall\textquotedblright{}, Old High German \emph{val} (EtymDB): it first meant siding with someone, then applause.
\end{cousinweb}

\begin{grammarlens}[title={What you've built}]
One more word for telling, explaining and describing at B1.
\end{grammarlens}

\subsection*{Your turn}
\begin{itemize}
\raggedright
  \item \emph{Say it:} \emph{der Beifall}
  \item \emph{Say it:} \emph{der Beifall}, once more
  \item \emph{Say it:} a sentence of your own with \emph{der Beifall}
  \item \emph{Recall:} say the German for the sparkling wine, then the German for to cheer, then say \emph{der Beifall} again
\end{itemize}

\begin{tcolorbox}[breakable,colback=teal!4,colframe=teal!35!black,title={Before you move on}]
What does der Beifall mean? (\textquotedblleft{}The applause\textquotedblright{}.) Say it once more.
\end{tcolorbox}
\section[applaudieren]{applaudieren — to applaud}

\label{lesson:GE-C308-applaudieren}



\begin{quote}
Before the new one: say the German for to cheer, then the German for the applause.
\end{quote}

\subsection*{You'll want to know: applaudieren}
\textbf{applaudieren} — \textquotedblleft{}to applaud\textquotedblright{}.

\begin{quote}
Die Leute applaudierten zehn Minuten lang.
\end{quote}

\begin{quote}
\emph{People applauded for ten minutes.}
\end{quote}

With the dative: \emph{Sie applaudierten dem Sänger}.

\begin{cousinweb}
\emph{applaudieren} is a learned borrowing from Latin \emph{applaudere} \textquotedblleft{}to strike upon, to applaud\textquotedblright{}, built on \emph{plaudere} (Lewis \& Short); EtymDB lists no parent.
\end{cousinweb}

\begin{grammarlens}[title={What you've built}]
One more word for telling, explaining and describing at B1.
\end{grammarlens}

\subsection*{Your turn}
\begin{itemize}
\raggedright
  \item \emph{Say it:} \emph{applaudieren}
  \item \emph{Say it:} \emph{applaudieren}, once more
  \item \emph{Say it:} a sentence of your own with \emph{applaudieren}
  \item \emph{Recall:} say the German for to cheer, then the German for the applause, then say \emph{applaudieren} again
\end{itemize}

\begin{tcolorbox}[breakable,colback=teal!4,colframe=teal!35!black,title={Before you move on}]
What does applaudieren mean? (\textquotedblleft{}To applaud\textquotedblright{}.) Say it once more.
\end{tcolorbox}
